\documentclass[11pt]{article}

\usepackage[T1]{fontenc}
\usepackage[utf8]{inputenc}
\usepackage{mathpazo}
\usepackage{microtype}
\usepackage[margin=1in,a4paper]{geometry}
\usepackage{amsmath,amssymb,amsthm}
\usepackage{enumitem}
\usepackage[hidelinks]{hyperref}
\usepackage{parskip}

\newtheorem{theorem}{Theorem}[section]
\newtheorem{proposition}[theorem]{Proposition}
\newtheorem{definition}[theorem]{Definition}
\newtheorem{corollary}[theorem]{Corollary}
\newtheorem{lemma}[theorem]{Lemma}
\newtheorem{remark}[theorem]{Remark}

\newcommand{\lf}[1]{\texttt{#1}}
\newcommand{\B}{\mathbb{B}}
\newcommand{\Occurs}{\mathrm{Occurs}}

\title{Undeniability Does Not Yield Succession}
\author{}
\date{}

\begin{document}
\maketitle
\vspace{-2em}

\begin{abstract}
A tokening that marks the denial of occurrence differently from its assertion
witnesses occurrence. That retorsive step is propositional once the tokening
fact is an antecedent, and it is formalizable over an arbitrary domain of loci.
Temporal succession is another matter. If claims are properties of
succession-free worlds, then two tensed structures over one world satisfy
exactly the same claims; in particular, a still structure and a structure with
nonempty succession are claim-indiscernible. Consequently, no undeniable claim
--- no claim true at every world that discriminates its denial --- can entail
running. Relative to the signature \((T,h,\mathrm{tok})\), an oriented
succession premise \(P\) is independent of every claim true at a shared
two-model witness: no such claim entails \(P\) for all tensed worlds, and none
refutes it. On loci with decidable equality, moreover, every
relabeling-invariant succession is symmetric, so orientation is not available
from invariance. Structural consequences of a bare distinction (two-valued
arena, uniqueness of complement, fixed-point freeness, period-two alternation)
are available as verified conditionals; none of them yields succession. The
results are checked in Lean~4 with an empty axiom footprint. The upshot is a
limit theorem for retorsive arguments: discrimination reaches occurrence; it
does not reach time.
\end{abstract}

\section{Introduction}
\label{sec:intro}

Retorsive and transcendental arguments aim to extract objective content from
the conditions of assertion or denial \cite{stroud1968,stern2000}. A standing
worry is that such arguments report only on what a denier must accept, leaving
the world untouched. The present paper isolates one place where the worry is
decisive and one place where it is not.

The positive half is familiar in outline. Let a medium classify propositions by
assigning each a locus and reading loci with a two-valued map. If the denial of
occurrence and its assertion receive different readings, then occurrence holds:
the two loci witness a distinction. Once the tokening difference is an
antecedent, the result is a propositional refutation of the conjunction of that
difference with the denial \cite{mackie1964,hintikka1962}. Section~\ref{sec:discrimination}
states the theorem.

The negative half is the contribution. Regiment claims as properties of
\emph{worlds} that carry loci, a reading, and a tokening, but no succession
relation. A tensed world adds succession. Under that typing, succession is
invisible to claims: any two tensed structures over one world satisfy exactly
the same claims (Section~\ref{sec:obstruction}). It follows that no claim true
wherever its denial is discriminated can entail running. Temporal realization
therefore requires a premise \(P\) outside the undeniable fragment. Relative to
the world signature, \(P\) is independent of every claim true at the two-model
witness of Section~\ref{sec:indep}. A further theorem locates the directional
content of \(P\): on loci with decidable equality, invariant succession is
always symmetric (Section~\ref{sec:orientation}).

The structural facts about a two-valued carrier --- that some equivariant law
moves a point exactly when the carrier has two elements, that the excluding
structureless law is complement, that complement has no fixed point and yields
a unique period-two orbit --- are recorded in Section~\ref{sec:structure} as
background, with proofs in the accompanying Lean development. They are not used
to derive succession. The paper's thesis is the limit, not an existence proof
for becoming.

Section~\ref{sec:dilemma} answers the objection that claims ought to be tensed.
Section~\ref{sec:related} situates the result. Section~\ref{sec:mech} records
the verification.

\section{Discrimination}
\label{sec:discrimination}

\begin{definition}[Occurrence]
\label{def:occurs}
For \(h : T \to \B\), write \(\Occurs(h)\) when there exist \(t_1,t_2 \in T\)
with \(h(t_1) \neq h(t_2)\).
\end{definition}

\begin{theorem}[Discrimination]
\label{thm:discrimination}
Let \(\mathrm{tok} : \mathrm{Prop} \to T\) and \(h : T \to \B\). If
\[
h(\mathrm{tok}(\lnot\Occurs(h))) \neq h(\mathrm{tok}(\Occurs(h))),
\]
then \(\Occurs(h)\) and the two tokens are distinct.
\end{theorem}

\begin{proof}
The two values of \(h\circ\mathrm{tok}\) are distinct elements of the image of
\(h\), so \(\Occurs(h)\) holds. If the tokens were equal, the values of \(h\)
at them would be equal.
\end{proof}

The argument is cardinality under a self-referential instance: any tokening
that takes two values shows that the medium has two elements; aiming the pair
at \(\lnot\Occurs\) and \(\Occurs\) refutes the denial. In Mackie's taxonomy
the result is absolute (propositional) once the tokening fact is an antecedent
\cite{mackie1964}. It is the explicit-tokening form of Hintikka's existential
inconsistency \cite{hintikka1962}.

The antecedent is not a theorem of the ambient type theory. A closed normal
term of type \(\mathrm{Prop}\to\B\) cannot depend on its argument, hence is
constant. Detachment is ostensive: some actual classification separates the
two readings. One concrete discharge uses the pair
\begin{quote}
Occurrence fails: every reading of this page classifies alike.\\
Occurrence holds: two readings of this page classify differently.
\end{quote}
Any \(h\) that separates those two readings satisfies the hypothesis. The
paper assumes such a discharge for the limit argument below; nothing later
requires more than the conditional. In particular, the obstruction does not
depend on which medium supplies the still discriminating world --- only on
there being one.

\begin{remark}
On \(\B\), distinctness coincides with being one complement apart. The third
conjunct of the verified form of Theorem~\ref{thm:discrimination} is therefore
redundant on two values. No temporal content is involved: a static pair of
marks already satisfies the theorem. The formal alternation of
Section~\ref{sec:structure} is likewise an \(\omega\)-indexed sequence of
values, a tenselessly specifiable object. Realization in succession is a
separate question, taken up next.
\end{remark}

\section{Worlds, freeze, and obstruction}
\label{sec:obstruction}

\begin{definition}[Worlds and claims]
\label{def:world}
A \emph{world} is a triple \(W=(T,h,\mathrm{tok})\): a type of loci, a reading
\(h:T\to\B\), and a tokening \(\mathrm{tok}:\mathrm{Prop}\to T\). A
\emph{claim} is a property of worlds. \(W\) \emph{discriminates the denial} of
a claim \(A\) when
\(h(\mathrm{tok}(\lnot A(W))) \neq h(\mathrm{tok}(A(W)))\). \(A\) is
\emph{undeniable} when \(A(W)\) holds at every world \(W\) that discriminates
the denial of \(A\).
\end{definition}

\begin{definition}[Tensed worlds]
\label{def:tensed}
A \emph{tensed world} is a world together with a binary succession relation on
its loci. It is \emph{still} when the relation is empty. It \emph{runs} when
there is an injective sequence of loci, each succeeding the last, whose
readings alternate under complement.
\end{definition}

Occurrence, read as a claim, is undeniable: Theorem~\ref{thm:discrimination}
transposed to worlds. Undeniability is nevertheless weak.

\begin{proposition}[Calibration]
\label{prop:calibration}
The constant-true claim is undeniable.
\end{proposition}

\begin{proof}
The consequent is trivial at every world.
\end{proof}

So undeniability constrains a claim only when the discrimination hypothesis
reaches the claim's content. Occurrence is such a claim; an arbitrary bridge to
running need not be.

\begin{proposition}[Freeze]
\label{prop:freeze}
For any world \(W\) and any two succession relations \(R_1,R_2\) on its loci,
every claim holds of the tensed world \((W,R_1)\) if and only if it holds of
\((W,R_2)\). In particular, over any inhabited world a still structure and a
structure with nonempty succession coexist and satisfy exactly the same claims.
\end{proposition}

\begin{proof}
Claims are typed over worlds, not over tensed worlds. The two tensed worlds
share \(W\), so they satisfy the same claims by definition. Equip \(W\) with
the empty relation and with the total relation for the second clause.
\end{proof}

\begin{theorem}[Obstruction]
\label{thm:obstruction}
Let \(A\) be undeniable and suppose \(A\) entails running. Then no still tensed
world discriminates the denial of \(A\).
\end{theorem}

\begin{proof}
Suppose a still tensed world discriminates the denial of \(A\). By
undeniability, \(A\) holds at its underlying world. By the entailment, the
tensed world runs. But a still world has no succeeding pair, so it does not
run.
\end{proof}

\begin{corollary}[No undeniable bridge]
\label{cor:nobridge}
If some still tensed world discriminates the denial of \(A\), then no
undeniable \(A\) entails running.
\end{corollary}

The corollary is a rebuttal schema against bridges to \emph{running}. Given a
candidate axiom offered as an undeniable entailment of \(\mathrm{Runs}\),
instantiate a still world that discriminates its denial (e.g.\ inscriptions on
a page with empty succession, once a separating reading is granted). The
candidate fails. The schema does not close the gap for candidates that are
never so tokened; it shows that undeniability, as defined, cannot be that
bridge. A parallel observation for the ordering half of \(P\) is recorded in
Section~\ref{sec:premise}.

Two readings of the loci are useful. On the \emph{act} reading, \(T\) comprises
readings performed by a reader; on the \emph{mark} reading, \(T\) comprises
inscriptions. The still witness for Corollary~\ref{cor:nobridge} uses marks
with empty succession. A reader who insists that the inscriptions were produced
in order is describing a second tensed structure over the same world; by
Proposition~\ref{prop:freeze} the two agree on every claim.

\section{The succession premise}
\label{sec:premise}

\begin{definition}[Oriented succession]
\label{def:P}
Fix a tensed world and a tokening. The \emph{oriented succession premise} at a
proposition \(X\) is
\[
P(X):\quad
\mathrm{succ}(\mathrm{tok}(\lnot X),\mathrm{tok}(X))
\;\land\;
\lnot\mathrm{succ}(\mathrm{tok}(X),\mathrm{tok}(\lnot X)).
\]
The first conjunct is the ordering half; the second is the orientation half.
\end{definition}

Over a fixed world, the empty succession makes every instance of \(P\) fail and
the total succession makes every ordering-half instance hold
(\lf{premise\_fails\_still}, \lf{premise\_holds\_running});
Proposition~\ref{prop:freeze} ensures the two structures agree on all claims.
Moreover, no undeniable claim can entail the ordering half of \(P\): if it did,
any still world discriminating that claim's denial would satisfy the claim and
hence witness a succession, contradicting stillness. (Corollary~\ref{cor:nobridge}
is the companion fact for entailments of \(\mathrm{Runs}\).) On the act
reading, a single instance of \(P\) is an ordinary report of reading order: the
denial was read, then the assertion, and not conversely \cite{moore1925}.
Warrant for one instance does not automatically extend to an infinite family of
instances.

\begin{theorem}[One instance]
\label{thm:actuality}
Given discrimination at \(X\) and one instance of \(P(X)\), the two tokens are
distinct, the successor's value is the complement of the predecessor's, and the
structure is not still.
\end{theorem}

\begin{proof}
Distinctness follows from discrimination (equal tokens yield equal readings).
The value claim is the two-element fact that distinct Boolean values are
complements. Non-stillness is immediate from the ordering conjunct.
\end{proof}

\begin{theorem}[Chain]
\label{thm:chain}
A chain of oriented, stepwise-discriminated instances over pairwise-distinct
loci yields an oriented run, hence a run, and refutes stillness.
\end{theorem}

\begin{proof}
Assemble the sequence of tokens. Injectivity is the pairwise-distinctness
hypothesis; each step contributes forward succession and failure of the
reverse; values alternate by the two-element fact as in
Theorem~\ref{thm:actuality}. Stillness fails at the first step.
\end{proof}

Theorem~\ref{thm:actuality} uses only the ordering conjunct of \(P\); the
orientation conjunct is idle in its proof. Orientation is needed for
distinctness without discrimination, and for oriented runs as opposed to mere
runs. Symmetric adjacency on \(\mathbb{N}\) --- the relation
\(n\,R\,m \leftrightarrow |n-m|=1\) --- runs and is not still, yet admits no
oriented run, since every edge is reversible.

Three commitments should be kept distinct:
\begin{enumerate}[label=(\arabic*),itemsep=1pt]
\item one oriented instance of \(P\) --- locally easy to grant on the act reading;
\item an infinite run --- a further assumption, not discharged by~(1);
\item A-theoretic passage --- a further gap beyond tenseless succession
\cite{williams1951,mctaggart1908}.
\end{enumerate}
The obstruction concerns the step from undeniability to~(1)--(2). Item~(3) is
left open even if~(2) is granted. The paper claims no derivation of~(2) or~(3)
from retorsion.

\section{Independence}
\label{sec:indep}

\begin{proposition}[Independence]
\label{prop:independence}
There is a world with two distinctly read loci such that:
\begin{enumerate}[label=(\roman*),itemsep=1pt]
\item the still tensed structure over it makes \(P\) fail for every tokening;
\item a tensed structure over it makes the ordering half of \(P\) hold for
every tokening;
\item every claim holds at both structures or at neither.
\end{enumerate}
Hence no claim \emph{true at that shared world} entails \(P\) for all tensed
worlds, and none such refutes \(P\)
(\lf{no\_claim\_settles\_premise}). The unrestricted reading is false: the
constant-false claim vacuously entails and refutes every target. The same
exhibition works for oriented \(P\) over a world whose tokening separates a
chosen pair. The full chain hypothesis is realized over a world with infinitely
many loci (total succession on \(\mathbb{N}\) with alternating reading, or
successor on \(\mathbb{N}\) for the oriented run).
\end{proposition}

\begin{proof}
Take \(T=\B\), \(h=\mathrm{id}\), and a constant tokening for the still/total
pair; Proposition~\ref{prop:freeze} gives~(iii), and the \(P\)-behaviour is as
above. For oriented independence, take succession to hold of \((\bot,\top)\)
alone. For the chain hypothesis, take \(T=\mathbb{N}\) with the total relation
and alternating reading (\lf{chain\_hypothesis\_realizable}); successor on
\(\mathbb{N}\) supplies the oriented run (\lf{oriented\_witness\_runs}).
\end{proof}

Independence is relative to the signature \((T,h,\mathrm{tok})\). If the
signature is enriched with causal facts and causal order determines temporal
order, \(P\) may become derivable \cite{mellor1998}. The present signature is
exactly what retorsion uses. Putnam's relativity argument concerns
frame-independent succession between spacelike-separated events
\cite{putnam1967}; successive readings by one reader are timelike.

\section{Orientation}
\label{sec:orientation}

Running constrains the supply of loci before it constrains succession: a
one- or two-element carrier admits no injective infinite sequence, whatever the
relation.

Call a succession relation \emph{invariant} when it is preserved by every
transposition of loci. (Loci are assumed to have decidable equality, so that
transpositions are definable from equality.)

\begin{proposition}[Classification]
\label{prop:invsucc}
An invariant succession is uniform on the diagonal and, given two distinct
loci, uniform off the diagonal: it is determined by the propositions
\(R(x,x)\) and \(R(x,y)\) for any \(x\neq y\). Equality and the empty relation
carry no run. An invariant succession that fails somewhere on the diagonal and
holds of some off-diagonal pair is disagreement. When the relation is
decidable (equivalently, \(\B\)-valued), the diagonal and off-diagonal values
each case into true or false, yielding exactly four extensionally possible
relations: empty, equality, disagreement, and total.
\end{proposition}

\begin{proof}[Proof sketch]
One transposition carries any diagonal pair to any other, so the diagonal is
uniform. Two transpositions carry any off-diagonal pair to any other when at
least two loci exist, so the off-diagonal is uniform. The verified content is
this biconditional uniformity (\lf{invariant\_succ\_classified}); a four-way
exhaustive classification for arbitrary \(\mathrm{Prop}\)-valued relations would
decide those two propositions and is not claimed. A succession contained in
equality identifies consecutive terms of any sequence, against injectivity.
The last claim: failure on the diagonal and success off it force
\(R(u,v)\leftrightarrow u\neq v\) by the two uniformities
(\lf{invariant\_says\_is\_ne}).
\end{proof}

\begin{theorem}[Invariant succession is symmetric]
\label{thm:orientation}
Every invariant succession is symmetric. Hence no invariant succession is
asymmetric on any pair, and no invariant succession carries an oriented run.
\end{theorem}

\begin{proof}
For \(u=v\) symmetry is trivial. For \(u\neq v\), off-diagonal uniformity
(Proposition~\ref{prop:invsucc}) yields \(R(u,v)\leftrightarrow R(v,u)\). An
oriented run requires a step whose reverse fails.
\end{proof}

Thus invariance determines the \emph{shape} of succession and not its
\emph{direction}. Three witnesses make the separation concrete:
\begin{enumerate}[itemsep=1pt,leftmargin=*]
\item symmetric adjacency on \(\mathbb{N}\) runs, is not still, and holds in
both directions at once;
\item the total relation satisfies the ordering half of \(P\) for every
tokening and the oriented half for none;
\item successor on \(\mathbb{N}\) yields an oriented run.
\end{enumerate}
A C-theorist who accepts order without direction \cite{farr2020,mctaggart1908}
is therefore stable under any invariance argument: symmetric adjacency is a
binary proxy for betweenness. The orientation clause of \(P\) is a genuine
addition to the ordering clause. What it secures is local irreversibility; a
succession may satisfy oriented-run and still contain a directed cycle
(successor on \(\mathbb{N}\) plus one reverse edge). Global topology remains
open.

\section{Background structure}
\label{sec:structure}

For completeness, the conditional structural results used only as background
are summarized here. Proofs are in \lf{lean/Becoming.lean}.

\begin{enumerate}[itemsep=2pt,leftmargin=*]
\item \emph{Arena.} On a carrier with decidable equality, there exists an
equivariant endofunction that moves some point if and only if the carrier has
exactly two elements (\lf{moves\_iff\_two}). For mark-respecting equivariance
(commutation with transpositions fixing a distinguished mark \(m\)), the
matching facts are: some mark-respecting map moves \(m\) iff the carrier has
two elements, and some mark-respecting map is fixed-point-free iff the carrier
has two elements (\lf{mark\_moves\_iff\_two}, \lf{marked\_refuses\_iff\_two}).
Moving an unmarked point is weaker: the constant map to \(m\) is
mark-respecting on any carrier and moves every \(x\neq m\)
(\lf{marked\_three\_moves}).
\item \emph{Uniqueness.} On \(\B\), the structureless (complement-equivariant)
unary laws are the identity and Boolean complement. Only complement excludes a
diagonal entry.
\item \emph{No fixed point; alternation.} Complement has no fixed point. Its
orbit from either seed is the unique period-two solution of the recursion; the
two orbits are relabelings. Fixed-point freeness is equivalent to stepwise
movement of every orbit.
\item \emph{Articulate relations.} Definable binary relations on \(\B\) that
exclude something and are satisfied (objectually or processually) are
extensionally inequality; satisfaction is then necessarily processual.
\end{enumerate}

Each item is a conditional or a classification. None mentions succession on
loci. Detaching occurrence for an actual medium uses the ostension of
Section~\ref{sec:discrimination}; detaching a temporal run uses \(P\).

\section{The regimentation dilemma}
\label{sec:dilemma}

The obstruction has the succession-blind typing of claims as an antecedent. An
objector may reply that tokenings are acts, acts are events, and the relevant
claims are therefore tensed \cite{apel1980,hintikka1962}. Grant the proposal.
Discrimination at the tensed level then either uses only the world components
or also uses the succession component.

\begin{theorem}[First horn]
\label{thm:hornone}
Suppose tensed-level discrimination depends on the world components alone. Let
\(A'\) hold at every tensed world whose world part discriminates, and let
\(A'\) entail running. Then no still tensed world discriminates.
\end{theorem}

\begin{proof}
As in Theorem~\ref{thm:obstruction}, one level up.
\end{proof}

\begin{theorem}[Second horn]
\label{thm:horntwo}
Suppose tensed-level discrimination requires that the two tokens stand in
succession. Then discrimination is an instance of the ordering half of \(P\),
and any predicate entailing a succession fact refutes stillness wherever it
holds.
\end{theorem}

\begin{proof}
The first claim is definitional. The second: a succeeding pair contradicts
stillness.
\end{proof}

On the first horn the obstruction survives. On the second, temporal content is
built into discrimination, so it is assumed rather than derived. Either way,
undeniability alone does not produce succession. The same boundary appears in
Kaplan's distinction between context-validity and necessity \cite{kaplan1989}:
a sentence may hold whenever tokened and still fail to be necessary.

\section{Related work}
\label{sec:related}

\paragraph{Transcendental arguments.}
Stroud's limit --- that such arguments reach the denier's commitments rather
than the world --- is the template \cite{stroud1968}.
Theorem~\ref{thm:discrimination} passes the limit for occurrence: an objective
distinction follows from a discrimination hypothesis.
Theorem~\ref{thm:obstruction} reinstates the limit for succession. In Stern's
terms the result is ambitious for occurrence and modest for succession
\cite{stern2000}. Illies's defence of ambitious transcendental arguments in
moral philosophy is among the targets of the obstruction if those arguments
require temporal or dynamical import from undeniability alone \cite{illies2003}.

\paragraph{Final justification.}
Albert's Münchhausen trilemma lists regress, circularity, and break-off
\cite{albert1985}. Apel and Kuhlmann seek to avoid break-off by performative
self-defeat \cite{apel1980,kuhlmann1985}; Habermas declines strong ultimate
grounding \cite{habermas1990}. The present position accepts break-off for the
temporal case: the conditionals are proved, \(P\) is independent of the claim
language, and the obstruction explains why an undeniable temporal bridge is
unavailable. The acceptance is local to succession; occurrence remains on the
ambitious side.

\paragraph{Hegel and Trendelenburg.}
Trendelenburg charged that Hegel's derivation of movement imports intuition
\cite{trendelenburg1870,hegel1832}. Theorem~\ref{thm:obstruction} is a formal
analogue under a succession-blind claim language: such a language entails no
succession. Section~\ref{sec:dilemma} defends the regimentation by showing that
the tensed alternative reproduces the same accounting. The paper does not offer
a reading of the Logic; it records that a presuppositionless combinatorial
beginning, if regimented as here, stops short of succession.

\paragraph{Temporal ontology.}
C-theory keeps order without direction \cite{mctaggart1908,farr2020}.
Theorem~\ref{thm:orientation} shows that refusal of direction is stable under
relabeling. The formal run, even when granted, is at-at variation along a
succession, not A-theoretic passage
\cite{williams1951,smart1949,prior1959,skow2015}. Mellor's token-reflexive
account matches the \(\mathrm{tok}/h\) signature and imposes the causal scope
condition of Section~\ref{sec:indep} \cite{mellor1998}. Grünbaum's
mind-dependent becoming places temporal content in awareness of events
\cite{grunbaum1967}, adjacent to a tokening-based account but silent on the
obstruction.

\paragraph{Laws of Form.}
Spencer-Brown's oscillating value \cite{spencerbrown1969} and related
three-valued or eigenform proposals \cite{varela1975,kauffman1987} address the
structural half of Section~\ref{sec:structure}. The present paper does not
re-derive that half; it separates it from succession and assigns the latter to
\(P\). Under mark-respecting equivariance, a third value does not yield
mark-moving dynamics (Section~\ref{sec:structure}, item~1).

\paragraph{Verified arguments.}
Machine-checked ontological arguments typically place existential content in
axioms and verify the derivation
\cite{oppenheimerzalta2011,benzmuller2014,benzmuller2016}. Here no axiom is
declared. Existential temporal content sits in \(P\) and in the ostensive
discharge of discrimination; the verified material is the surrounding
conditional and limit structure.

\section{Verification}
\label{sec:mech}

The definitions and theorems of Sections~\ref{sec:discrimination}--\ref{sec:dilemma}
are named declarations in two standalone Lean~4 files \cite{demoura2021},
prelude only, with no imports and no axioms. Each file ends with a
\lf{\#print axioms} audit; every line reports independence from all axioms,
including Lean's standard three (\(\mathrm{propext}\),
\(\mathrm{Classical.choice}\), \(\mathrm{Quot.sound}\)).

\begin{center}
\begin{tabular}{@{}ll@{}}
\hline
Result & Lean \\
\hline
Theorem~\ref{thm:discrimination} & \lf{event\_retorsion}, \lf{discrimination} \\
Definition~\ref{def:world} & \lf{World}, \lf{Claim}, \lf{Undeniable}, \ldots \\
Proposition~\ref{prop:calibration} & \lf{undeniable\_trivial} \\
Proposition~\ref{prop:freeze} & \lf{freeze}, \lf{two\_models} \\
Theorem~\ref{thm:obstruction} & \lf{obstruction}, \lf{still\_never\_runs} \\
Corollary~\ref{cor:nobridge} & \lf{no\_undeniable\_bridge} \\
Definition~\ref{def:P} & \lf{Successive}, \lf{OrientedSuccessive} \\
Theorem~\ref{thm:actuality} & \lf{oriented\_tick} \\
Theorem~\ref{thm:chain} & \lf{oriented\_chain\_runs} \\
Proposition~\ref{prop:independence} & \lf{premise\_independent},
  \lf{premise\_fails\_still}, \lf{premise\_holds\_running},
  \lf{no\_claim\_settles\_premise}, \lf{chain\_hypothesis\_realizable},
  \lf{oriented\_witness\_runs} \\
Proposition~\ref{prop:invsucc} & \lf{invariant\_succ\_classified},
  \lf{invariant\_says\_is\_ne} \\
Theorem~\ref{thm:orientation} & \lf{invariant\_succ\_symm}, \lf{no\_invariant\_orientation} \\
Theorems~\ref{thm:hornone},~\ref{thm:horntwo} & \lf{obstruction\_tensed}, \lf{horn\_two} \\
Section~\ref{sec:structure} & \lf{lean/Becoming.lean}
  (\lf{moves\_iff\_two}, \lf{mark\_moves\_iff\_two}, \ldots) \\
\hline
\end{tabular}
\end{center}

\vspace{0.5em}
\noindent
Check with \lf{lean lean/Becoming.lean} and \lf{lean lean/Obstruction.lean}
against toolchain \lf{leanprover/lean4:v4.31.0}.

\section{Conclusion}
\label{sec:conclusion}

Discrimination yields occurrence. Under a succession-blind claim language, it
does not yield succession: undeniability and running are incompatible once a
still world that discriminates the relevant denial is granted --- freeze
supplies the still structure over a given world, and discrimination at that
world is the ostensive step of Section~\ref{sec:discrimination}. The succession
premise \(P\) is independent of every claim true at the two-model witness,
relative to the retorsive signature, and its orientation clause is invisible to
invariant structure. The a priori core is therefore the conditionals together
with this limit. Time, if wanted, is paid for separately.

\begin{thebibliography}{99}

\bibitem{albert1985} H.\ Albert, \emph{Treatise on Critical Reason}, trans.\
M.\ V.\ Rorty, Princeton University Press, 1985.

\bibitem{apel1980} K.-O.\ Apel, \emph{Towards a Transformation of Philosophy},
trans.\ G.\ Adey and D.\ Frisby, Routledge \& Kegan Paul, 1980.

\bibitem{benzmuller2014} C.\ Benzm\"uller and B.\ Woltzenlogel Paleo,
``Automating G\"odel's Ontological Proof of God's Existence with Higher-order
Automated Theorem Provers,'' in \emph{ECAI 2014}, IOS Press, 2014, 93--98.

\bibitem{benzmuller2016} C.\ Benzm\"uller and B.\ Woltzenlogel Paleo, ``The
Inconsistency in G\"odel's Ontological Argument,'' in \emph{Proceedings of
IJCAI 2016}, 2016, 936--942.

\bibitem{demoura2021} L.\ de Moura and S.\ Ullrich, ``The Lean 4 Theorem Prover
and Programming Language,'' in \emph{Automated Deduction: CADE 28}, Springer,
2021, 625--635.

\bibitem{farr2020} M.\ Farr, ``C-Theories of Time: On the Adirectionality of
Time,'' \emph{Philosophy Compass} 15(12), 2020.

\bibitem{grunbaum1967} A.\ Gr\"unbaum, ``The Status of Temporal Becoming,'' in
\emph{Modern Science and Zeno's Paradoxes}, Wesleyan University Press, 1967,
7--36.

\bibitem{habermas1990} J.\ Habermas, \emph{Moral Consciousness and
Communicative Action}, trans.\ C.\ Lenhardt and S.\ W.\ Nicholsen, MIT Press,
1990.

\bibitem{hegel1832} G.\ W.\ F.\ Hegel, \emph{The Science of Logic}, trans.\
G.\ di Giovanni, Cambridge University Press, 2010.

\bibitem{hintikka1962} J.\ Hintikka, ``Cogito, Ergo Sum: Inference or
Performance?,'' \emph{Philosophical Review} 71(1), 1962, 3--32.

\bibitem{illies2003} C.\ Illies, \emph{The Grounds of Ethical Judgement},
Oxford University Press, 2003.

\bibitem{kaplan1989} D.\ Kaplan, ``Demonstratives,'' in \emph{Themes from
Kaplan}, ed.\ J.\ Almog, J.\ Perry, and H.\ Wettstein, Oxford University Press,
1989, 481--563.

\bibitem{kauffman1987} L.\ H.\ Kauffman, ``Self-Reference and Recursive Forms,''
\emph{Journal of Social and Biological Structures} 10(1), 1987, 53--72.

\bibitem{kuhlmann1985} W.\ Kuhlmann, \emph{Reflexive Letztbegr\"undung}, Karl
Alber, 1985.

\bibitem{mackie1964} J.\ L.\ Mackie, ``Self-Refutation: A Formal Analysis,''
\emph{Philosophical Quarterly} 14(56), 1964, 193--203.

\bibitem{mctaggart1908} J.\ M.\ E.\ McTaggart, ``The Unreality of Time,''
\emph{Mind} 17(68), 1908, 457--474.

\bibitem{mellor1998} D.\ H.\ Mellor, \emph{Real Time II}, Routledge, 1998.

\bibitem{moore1925} G.\ E.\ Moore, ``A Defence of Common Sense,'' in
\emph{Contemporary British Philosophy, Second Series}, ed.\ J.\ H.\ Muirhead,
George Allen \& Unwin, 1925, 192--233.

\bibitem{oppenheimerzalta2011} P.\ E.\ Oppenheimer and E.\ N.\ Zalta, ``A
Computationally-Discovered Simplification of the Ontological Argument,''
\emph{Australasian Journal of Philosophy} 89(2), 2011, 333--349.

\bibitem{prior1959} A.\ N.\ Prior, ``Thank Goodness That's Over,''
\emph{Philosophy} 34(128), 1959, 12--17.

\bibitem{putnam1967} H.\ Putnam, ``Time and Physical Geometry,'' \emph{Journal
of Philosophy} 64(8), 1967, 240--247.

\bibitem{skow2015} B.\ Skow, \emph{Objective Becoming}, Oxford University Press,
2015.

\bibitem{smart1949} J.\ J.\ C.\ Smart, ``The River of Time,'' \emph{Mind}
58(232), 1949, 483--494.

\bibitem{spencerbrown1969} G.\ Spencer-Brown, \emph{Laws of Form}, George Allen
and Unwin, 1969.

\bibitem{stern2000} R.\ Stern, \emph{Transcendental Arguments and Scepticism},
Oxford University Press, 2000.

\bibitem{stroud1968} B.\ Stroud, ``Transcendental Arguments,'' \emph{Journal of
Philosophy} 65(9), 1968, 241--256.

\bibitem{trendelenburg1870} F.\ A.\ Trendelenburg, \emph{Logische
Untersuchungen}, vol.\ 1, 3rd ed., S.\ Hirzel, Leipzig, 1870.

\bibitem{varela1975} F.\ J.\ Varela, ``A Calculus for Self-Reference,''
\emph{International Journal of General Systems} 2(1), 1975, 5--24.

\bibitem{williams1951} D.\ C.\ Williams, ``The Myth of Passage,'' \emph{Journal
of Philosophy} 48(15), 1951, 457--472.

\end{thebibliography}

\end{document}
